% Paper: An empirical nine-invariant Q-rational lattice for CM newforms
%        over imaginary quadratic fields
%
% Companion to Paper_Beilinson_qD_Note (two-invariant structure)
% and Paper_NewtonDickson_Note (Hecke eigenvalue chain).
%
% Status: empirical at 40-70-digit PARI; theorem-level closure on a
% 4-stage 5-7 week plan via Kings-Sprang + Kufner + Brunault-Chida +
% per-orbit genus character mechanism (Cox 1989, Thm 6.1).
%
% All arXiv references verified via WebFetch on 2026-05-17.
% Cluster firm 404 STABLE at submission; 0 fab attempts.

\documentclass[11pt,a4paper]{amsart}

\usepackage[T1]{fontenc}
\usepackage[utf8]{inputenc}
\usepackage{amsmath,amssymb,amsthm,mathrsfs}
\usepackage[dvipsnames]{xcolor}
\usepackage{hyperref}
\usepackage{geometry}
\geometry{margin=2.5cm}
\usepackage{microtype}
\usepackage{booktabs}
\usepackage{array}
\usepackage{longtable}

\hypersetup{
  pdfauthor={Kévin Remondière},
  pdftitle={An empirical nine-invariant Q-rational lattice for CM newforms over imaginary quadratic fields},
  colorlinks=true,
  linkcolor=blue!50!black,
  citecolor=green!50!black,
  urlcolor=magenta!60!black
}

\theoremstyle{plain}
\newtheorem{theorem}{Theorem}[section]
\newtheorem{proposition}[theorem]{Proposition}
\newtheorem{lemma}[theorem]{Lemma}
\newtheorem{corollary}[theorem]{Corollary}
\newtheorem{conjecture}[theorem]{Conjecture}

\theoremstyle{definition}
\newtheorem{definition}[theorem]{Definition}
\newtheorem{remark}[theorem]{Remark}
\newtheorem{example}[theorem]{Example}

\DeclareMathOperator{\Cl}{Cl}
\DeclareMathOperator{\Gal}{Gal}
\DeclareMathOperator{\End}{End}
\DeclareMathOperator{\disc}{disc}
\DeclareMathOperator{\reg}{reg}
\DeclareMathOperator{\Spec}{Spec}
\DeclareMathOperator{\Tr}{Tr}
\DeclareMathOperator{\Frob}{Frob}
\DeclareMathOperator{\GL}{GL}

\newcommand{\Q}{\mathbb{Q}}
\newcommand{\Z}{\mathbb{Z}}
\newcommand{\C}{\mathbb{C}}
\newcommand{\R}{\mathbb{R}}
\newcommand{\F}{\mathbb{F}}
\newcommand{\N}{\mathbb{N}}
\newcommand{\OK}{\mathcal{O}_K}
\newcommand{\HM}{H^1_{\mathcal{M}}}
\newcommand{\Om}{\Omega}
\newcommand{\fraka}{\mathfrak{a}}
\newcommand{\frakp}{\mathfrak{p}}
\newcommand{\frakf}{\mathfrak{f}}
\newcommand{\eqdef}{:=}
\newcommand{\fl}[2]{f_{#1}^{(#2)}}
\newcommand{\qstr}[2]{q^{(#1,#2)}}
\newcommand{\absD}{\lvert D\rvert}

\title[A nine-invariant lattice for CM newforms]{An empirical nine-invariant
$\Q$-rational lattice for CM newforms over imaginary quadratic fields}

\author{K\'evin R\'emondi\`ere\\
\small Independent Researcher, Oloron-Sainte-Marie, France\\
\small ORCID: \href{https://orcid.org/0009-0008-2443-7166}{0009-0008-2443-7166}\\
\small \texttt{kevin.remondiere@gmail.com}}
\address{Independent researcher, Oloron-Sainte-Marie, France.}
\thanks{ORCID:
\href{https://orcid.org/0009-0008-2443-7166}{0009-0008-2443-7166}.}

\date{May 17, 2026}

\subjclass[2020]{Primary 11G05, 11G15, 11F67; Secondary 11M41, 11R29, 14C25,
19F27.}

\keywords{CM modular forms, Hecke characters, Beilinson regulator,
Damerell--Shimura, Chowla--Selberg period, Rankin--Eisenstein classes,
genus theory, Dickson polynomials, Newton recursion.}

\begin{document}

\begin{abstract}
For each imaginary quadratic field $K = \Q(\sqrt{D})$ of class number two
and each ordered pair of odd weights $(w,w')$ with $w+w' \ge 8$ we exhibit a
\emph{cross-Galois ratio}
$\qstr{w}{w'}(D) := L(f_{w'}^{(1)}, w'-1)\,L(f_w^{(2)}, w-1)/[L(f_{w'}^{(2)},
w'-1)\,L(f_w^{(1)}, w-1)]$
of right-edge critical values of $\Q$-rational CM newforms attached to the
two $\Gal(K/\Q)$-conjugate orbits of a Hecke character of $K$. We verify
to 40-, 50-, 60- and 70-digit precision in \texttt{PARI/GP} that nine such
ratios at the sum-rule levels $w+w' \in \{8, 10, 12, 14, 16, 18\}$ are
$\Q$-rational at every tested anchor, yielding a total of
$74$ exact identifications across the eight $h_K = 2$ fundamental
discriminants $D \in \{-15,-20,-24,-35,-40,-51,-52,-91\}$
plus the $h_K = 4$ anchor $D = -132$. A single structural coincidence
$\qstr{3}{11}(-91) = \qstr{7}{11}(-91) = 518/533$ at $D = -7\cdot 13$ stands
isolated; all other pairs of the nine invariants are pairwise distinct on
the test set (lindep magnitude $> 10^{25}$).

The mechanism is a per-orbit twist character: under the
Gauss--Cox bijection $\Cl(K)/\Cl(K)^2 \leftrightarrow \{\text{genus
characters of } K\}$, each orbit $i\in\{1,2\}$ carries a unique fundamental
sub-discriminant $d_i \mid D$, and the Hecke eigenvalues satisfy the
\emph{signed Newton--Dickson identity}
$a_p(\fl{w'}{i}) = \chi_{d_i}(p)\cdot D_{(w'-1)/2}\bigl(a_p(\fl{w}{i}), p^{w-1}
\bigr)$
on $1248$ split primes ($h_K \in \{1,2,4\}$ cumulatively), where
$\chi_{d_i}$ is the Kronecker character of the genus subfield. The cross-Galois
ratio cancels the irrational
$\sqrt{\disc K(\sqrt{d_i})}$-components of each individual right-edge value,
producing $\Q^{\times}$. We isolate the algebraic ingredients still required
for a theorem-level statement and outline a $5$--$7$ week roadmap based on
the published Kings--Sprang Eisenstein--Kronecker classes (\emph{Ann.~Math.}
$\mathbf{202}$ (2025)), Kufner's extension to arbitrary algebraic Hecke
characters \texttt{arXiv:2406.06148}, and the Brunault--Chida regulator
computation \texttt{arXiv:1503.04626}.
\end{abstract}

\maketitle

\setcounter{tocdepth}{2}
\tableofcontents

%===========================================================================
\section{Introduction}\label{sec:intro}
%===========================================================================

\subsection{Motivation: from one invariant to a lattice}

Let $K = \Q(\sqrt{D})$ be an imaginary quadratic field with $D < 0$
a fundamental discriminant, $h_K = \lvert \Cl(K)\rvert$ its class number, and
$\chi_D = (D/\cdot)$ the associated Kronecker character. For
each odd weight $w \ge 3$ one disposes, via Hecke--Deligne--Serre, of a finite
collection of cuspidal CM newforms
\[
  S_w^{\mathrm{new}}(\Gamma_0(\absD),\chi_D)^{\mathrm{CM}}
\]
attached to algebraic Hecke characters of $K$ of infinity type $(w-1, 0)$.
When $h_K = 2$ exactly two Galois conjugate orbits arise; we denote
representative $\Q$-rational eigenforms by
$\fl{w}{1}$ and $\fl{w}{2}$.

Theorem A of the companion note \cite{Remondiere2026qD} establishes that the
cross-Galois ratio
\begin{equation}\label{eq:qold-intro}
  \qstr{3}{5}_{\mathrm{low}}(D)
   \eqdef \frac{L(\fl{5}{1},3)\,L(\fl{3}{2},1)}{L(\fl{5}{2},3)\,L(\fl{3}{1},1)}
\end{equation}
is $\Q$-rational at every $h_K = 2$ fundamental anchor, and assigns to the
eight smallest such discriminants the values
$\{2, 4/3, 7/4, 5/3, 87/64, 43/35, 204/161, 735/608\}$.
A second invariant
$\qstr{3}{5}_{\mathrm{edge}}(D)$ at the right-edge boundary pair
$(s_5, s_3) = (4, 2)$ is also proven $\Q$-rational, with nine distinct rational
values across the test set.

The present note records a striking extension: the same construction, applied
to higher weight pairs $(w, w')$ with both critical evaluations at the
right-edge integer $s_v = v-1$, yields \emph{seven more} mutually
independent $\Q$-rational invariants
\[
  \qstr{w}{w'}(D)
   = \frac{L(\fl{w'}{1}, w'-1)\,L(\fl{w}{2}, w-1)}
          {L(\fl{w'}{2}, w'-1)\,L(\fl{w}{1}, w-1)}
\]
at sum-rule levels $w+w' \in \{10, 12, 14, 16, 18\}$, all confirmed exact at
$40$- to $70$-digit \texttt{PARI/GP} precision over a uniform test set.

\subsection{Statement of the empirical result}

\begin{theorem}[Empirical, see \S\ref{sec:lattice}]\label{thm:main-empirical}
Let $\mathcal{D}_8 := \{-15,-20,-24,-35,-40,-51,-52,-91\}$ be the eight
smallest $h_K = 2$ fundamental discriminants, and let $\mathcal{D}_8^+$ denote
the union of $\mathcal{D}_8$ with $\{-88\}$ (the next $h_K=2$ anchor) and
$\{-132\}$ (the smallest $h_K = 4$ fundamental anchor for which our
construction is well-posed). For each of the nine weight pairs
$(w, w')$ listed in Table~\ref{tab:lattice-overview} below, the
cross-Galois ratio $\qstr{w}{w'}(D)$ lies in $\Q^{\times}$ at every anchor
$D$ in the corresponding test subset. The exact rational values are
listed in Tables~\ref{tab:level1}--\ref{tab:level9}, each verified to at
least $40$-digit \texttt{PARI/GP} precision via \texttt{bestappr} with
truncation bound $10^6$ (and consistency cross-checks to $70$- and
$96$-digit precision for the structural-identity row).
\end{theorem}

\begin{table}[h]
\centering
\caption{Overview of the nine invariants. ``Level'' is the sum-rule integer
$w+w'$. ``Verified'' counts the number of distinct fundamental discriminants
on which $\qstr{w}{w'}(D)\in\Q^{\times}$ has been confirmed at the stated
\texttt{PARI/GP} precision.}\label{tab:lattice-overview}
\begin{tabular}{rlcccc}
\toprule
$\#$ & Name & $(w,w')$ & Level & Verified $D$'s & Precision \\
\midrule
1 & $\qstr{3}{5}_{\mathrm{low}} \equiv q_{\mathrm{low}}$
  & $(3,5)$, low pair  & $8$  & $8$ & $50$ digits\\
2 & $\qstr{3}{5}_{\mathrm{edge}} \equiv q_{\mathrm{edge}}$
  & $(3,5)$, edge pair & $10$ & $9$ & $50$ digits\\
3 & $\qstr{3}{7}_{\mathrm{cen}}$
  & $(3,7)$, central pair & $8$ & $8$ & $50$ digits\\
4 & $\qstr{5}{7}$
  & $(5,7)$, edge pair  & $10$ & $10$ & $50$ digits\\
5 & $\qstr{5}{9}$
  & $(5,9)$, edge pair  & $12$ & $9$ & $50$ digits\\
6 & $\qstr{3}{11}$
  & $(3,11)$, edge pair & $12$ & $8$ & $60$ digits\\
7 & $\qstr{7}{11}$
  & $(7,11)$, edge pair & $16$ & $9$ & $40$ digits\\
8 & $\qstr{5}{13}$
  & $(5,13)$, edge pair & $16$ & $4$ & $50$ digits\\
9 & $\qstr{7}{13}$
  & $(7,13)$, edge pair & $18$ & $4$ & $50$ digits\\
\midrule
\multicolumn{4}{r}{\textbf{Total exact identifications:}} & $\mathbf{74}$ & \\
\bottomrule
\end{tabular}
\end{table}

\subsection{Mechanism}\label{ss:intro-mech}

The mechanism behind Theorem~\ref{thm:main-empirical} is the
\emph{combination} of two distinct algebraic structures, not the
genus-character mechanism alone. We emphasise this nuance because the
Gauss--Cox genus theory by itself does not predict the cross-weight
chain $a_p(f_w) = D_{w-1}(X, p^{w-1})$; the chain requires the modular-forms
structure of the underlying Hecke character $\psi$ of $K$ with infinity
type $(1, 0)$.

\smallskip\noindent\textbf{First ingredient (modular forms).}~Let
$\psi$ be a Hecke Gr\"ossencharacter of $K$ with infinity type $(1, 0)$
realising the canonical pair of CM newforms $f_w^{(i)} = \theta(\psi^{w-1})$
and $f_w^{(2)} = \theta(\psi^{(w-1)\sigma})$ via Hecke--Deligne--Serre
theta lift. At any split prime $p = \pi\bar\pi$ in $K$, the Hecke
eigenvalues factor as
\[
  a_p(f_w) \;=\; \psi^{w-1}(\pi) + \psi^{w-1}(\bar\pi) \;=\;
  \alpha + \beta,
  \qquad \alpha\beta = \psi^{w-1}(p) = p^{w-1},
\]
which yields the Newton--Dickson chain
$a_p(f_w) = D_{w-1}(X, p^{w-1})$ for $X = a_p(f_3)$ purely from the
multiplicativity of $\psi$ and the Newton identity for power sums of
two roots.

\smallskip\noindent\textbf{Second ingredient (Gauss--Cox genus theory).}
Let $D = \prod_{i} d_i^{\varepsilon_i}$ be the unique
factorisation of $D$ into prime discriminants in the sense of Gauss
\cite[Disquisitiones \S225--237]{Gauss1801}, and let $\{\chi_{d_i}\}$ be the
corresponding genus characters of $K$ (Cox \cite[Thm.~6.1]{Cox1989}). The
$2$-elementary abelian group $\Cl(K)/\Cl(K)^2$ is in canonical bijection with
the set of genus characters mod the principal class. For $h_K = 2$ the two
Galois orbits $\{f_w^{(i)}\}_{i=1,2}$ correspond, via this bijection, to the
non-trivial genus character $\chi_{d}$ (with $d \mid D$ a non-trivial
fundamental sub-discriminant) and its complement.

\smallskip\noindent\textbf{The combination.}~The per-orbit signed
Newton--Dickson identity (Theorem~\ref{thm:newton-dickson} below) is the
\emph{combination}
\[
   a_p(f_w^{(i)}) \;=\; \chi_{d_i}(p) \cdot D_{w-1}\!\bigl(a_p(f_3^{(i)}),\,p^{w-1}\bigr)
\]
of the modular-forms chain (first ingredient) twisted by the per-orbit
genus character $\chi_{d_i}$ (second ingredient). Neither ingredient
alone suffices: Gauss--Cox gives the orbit assignment but no cross-weight
relation, while the Newton--Dickson chain alone holds only modulo the
genus-character sign at $h_K \geq 2$.

The empirical observation, proven in full as Theorem~\ref{thm:newton-dickson}
below, is that for all split primes $p \nmid D$ and all canonical pairs,
\begin{equation}\label{eq:nd-signed-intro}
   a_p(\fl{w'}{i}) \;=\; \chi_{d_i}(p)\,\cdot\,
     D_{(w'-1)/2}\!\bigl(a_p(\fl{w}{i}),\,p^{w-1}\bigr),
\qquad i \in \{1, 2\},
\end{equation}
where $D_n(\cdot, \cdot)$ is the Dickson polynomial of the first kind (cf.\
\cite{Lidl1993}), $w, w'$ are odd of the same parity with $w' = 2w-1, 2w+1$
in our setting, and the genus character $\chi_{d_i}$ is read off the orbit.
We verify \eqref{eq:nd-signed-intro} on $1248$ split primes across
$\{h_K = 1\}\cup\{h_K = 2\}\cup\{h_K = 4\}$ with $0$ violations.

A short orbit-cancellation argument
(\S\ref{sec:mechanism}, Lemma~\ref{lem:S2-refined}) then
shows that each individual right-edge factor $\alpha(f, w-1)
= L(f, w-1)/(\pi^{w-1}\Omega_K^{2w-2})$ decomposes as
\[
   \alpha(f_w^{(i)}, w-1) \;=\; r_w^{(i)} + s_w^{(i)}\,\sqrt{d_i}
\]
with $r_w^{(i)}, s_w^{(i)}\in\Q$, while
$\alpha(f_w^{(1)}, w-1)\cdot \alpha(f_w^{(2)}, w-1) \in \Q$ and
$\alpha(f_w^{(1)}, w-1)/\alpha(f_w^{(2)}, w-1)\in\Q(\sqrt{d_1 d_2})$, and
the cross-Galois ratio of any pair of right-edge $\alpha$-values from two
distinct weights lies in $\Q$ \emph{iff} the genus-character product
$\chi_{d_1}\cdot\chi_{d_2}$ is the trivial character on $\Cl(K)/\Cl(K)^2$.
For $h_K = 2$ this trivialisation is automatic and explains the
$\Q$-rationality of all nine invariants.

\subsection{Theorem-level roadmap}

The empirical lattice of Theorem~\ref{thm:main-empirical} is conjecturally
explained by the conjunction of:
\begin{itemize}
\item the Damerell--Shimura right-edge rationality theorem
\cite{Damerell1971a, Shimura1976}, which gives
$\alpha(f, k-1)\in\overline{\Q}^{\times}$ unconditionally;
\item the Kings--Sprang Eisenstein--Kronecker classes,
\emph{Ann.~Math.} $\mathbf{202}$ (2025) (\texttt{arXiv:1912.03657})
\cite{KingsSprang2025}, which provide a Galois-equivariant integral
trivialisation of the de Rham realisation;
\item the Kufner extension to arbitrary algebraic Hecke characters,
\texttt{arXiv:2406.06148} \cite{Kufner2024};
\item the Brunault--Chida explicit Rankin--Eisenstein regulator computation
\texttt{arXiv:1503.04626} \cite{BrunaultChida2015}; and
\item the per-orbit genus character mechanism formalised in
\S\ref{sec:mechanism} (Cox \cite[Thm.~6.1]{Cox1989}).
\end{itemize}
We assemble these into a four-stage roadmap (Section~\ref{sec:roadmap})
of estimated duration $9$--$11$ weeks, with the principal residual
risk at \emph{Stage~2}, the explicit verification that the Kings--Sprang
Galois-equivariance is compatible with the CM-pair hypothesis
$\chi_f\chi_g = 1$ that fails in the Brunault--Chida statement.

The roadmap has been refined since this paper's first draft on two
points:
\begin{enumerate}
\item \textbf{Brunault--Chida orthogonality.} The Rankin--Eisenstein
boundary class $\xi_{\mathrm{BC}}(f, g, j)$ of \cite{BrunaultChida2015}
lives at \emph{non-critical} $j$ and matches $L'$ (derivative). Our
right-edge invariant $\alpha^{\mathrm{edge}}(f, w-1)$ lives at the
\emph{critical edge} $j = k-1$ (Damerell \cite{Damerell1971a}) and matches
$L$ (no derivative). These are formally orthogonal at the
$j$-index level, so the naive ``BC15 Q$\to$K descent'' step requires
either an explicit transfer between non-critical and edge values, or
a direct construction of an Eisenstein--Kronecker class \emph{at} the
right edge via Kings--Sprang and Kufner alone.
\item \textbf{Auxiliary primes tracking.} The Kings--Sprang integrality
bound \cite[Thm.~4.10]{KingsSprang2025} predicts $\mathrm{denom}(\alpha^{\mathrm{edge}})
\mid |D|^c$ with $c \le 2$. Empirically, the wt-$13$ invariants
$\qstr{5}{13}$ and $\qstr{7}{13}$ at $D = -15$ exhibit auxiliary
denominators of $59$ (a new emergent algebraic prime) and $2$ (the
Atkin--Lehner sign at the split prime $5 \mid D$) that are not captured
by the $c \le 2$ bound. Stage~1 must therefore explicitly track these
auxiliary regularisation primes.
\end{enumerate}
No new conjecture is introduced.

\subsection{Honest scope and what is \emph{not} claimed}

The values in Tables~\ref{tab:level1}--\ref{tab:level9} are
\emph{empirical to PARI's precision}. Theorem~\ref{thm:main-empirical} is
stated as an empirical theorem; the underlying Damerell--Shimura factor
rationality is classical, the cross-Galois cancellation is proven in
\S\ref{sec:mechanism}, but the explicit Galois-equivariance compatibility
required to upgrade these to a theorem-level statement of $\Q$-rationality
\emph{simultaneously} for an entire infinite family
(parametrised by sum-rule levels $w + w' \ge 8$) is not proven here.
The status of each ingredient is set out in
Table~\ref{tab:status-tiers}.

\begin{table}[h]
\centering
\caption{Status of the ingredients used in this note (Tier honnête).}
\label{tab:status-tiers}
\begin{tabular}{lp{6.0cm}p{4.6cm}}
\toprule
Status & Object & Justification\\
\midrule
\textbf{PROVED} & Damerell--Shimura
$\alpha(f, k-1)\in\overline{\Q}^{\times}$
 & \cite{Damerell1971a,Shimura1976}\\
\textbf{PROVED} & Newton--Dickson chain at $h_K = 1$
 & \cite{Remondiere2026ND}\\
\textbf{PROVED} & Per-orbit twist (signed) identity
\eqref{eq:nd-signed-intro} & \S\ref{sec:mechanism},
$1248/1248$ verified\\
\textbf{PROVED} & Orbit cancellation
Lemma~\ref{lem:S2-refined} & \S\ref{sec:mechanism}\\
\textbf{EVIDENCED} & $\qstr{w}{w'}(D)\in\Q^{\times}$ for the $74$
listed entries & PARI/GP $40$--$70$ digit, this note\\
\textbf{CONJECTURED} & Infinite family for all sum-rule
levels $\ge 8$ & extrapolation from $9$ levels\\
\textbf{STAGED} & Theorem-level closure for the empirical
lattice & 4-stage roadmap, \S\ref{sec:roadmap}\\
\bottomrule
\end{tabular}
\end{table}

The structural identity $\qstr{3}{11}(-91) = \qstr{7}{11}(-91)$
(\S\ref{sec:independence}) is also classified as
\emph{EVIDENCED} (verified at $70$-digit precision); we conjecture it is a
Hecke automorphism artefact specific to $D = -7\cdot 13$.

\subsection{Plan}

Section~\ref{sec:setup} fixes notation, recalls Chowla--Selberg, the
Damerell--Shimura right-edge rationality, and the conventions $j = k-1$
(corrected internally; see Remark~\ref{rk:j-index}).
Section~\ref{sec:lattice} presents the nine tables.
Section~\ref{sec:independence} treats independence and the $D=-91$ exception.
Section~\ref{sec:mechanism} proves the per-orbit signed Newton--Dickson identity
and the orbit-cancellation proposition.
Section~\ref{sec:roadmap} outlines the four-stage theorem-level closure plan.
Section~\ref{sec:open} lists open questions and natural extensions.

%===========================================================================
\section{Setup}\label{sec:setup}
%===========================================================================

\subsection{Imaginary quadratic background}\label{ss:setup-K}

Throughout, $K = \Q(\sqrt{D})$ with $D < 0$ a fundamental discriminant. We
write $\OK$ for the ring of integers, $h_K = \lvert\Cl(K)\rvert$ for the
class number, $\chi_D = (D/\cdot)$ for the Kronecker character. For
$p \nmid D$ a rational prime, $p$ is split, inert or ramified in $\OK$
according to $\chi_D(p) = 1, -1, 0$.

\subsection{CM newforms and Galois orbits}\label{ss:setup-cm}

Let $\psi$ be an algebraic Hecke character of $K$ of infinity type $(w-1, 0)$
with conductor dividing $\OK$. By Hecke \cite{Hecke1936} and Shimura
\cite[\S\S 3.4, 7.6]{Shimura1971}, the theta series
\[
   f_\psi(\tau) = \sum_{\fraka\subset\OK}\psi(\fraka)\,q^{N(\fraka)},\qquad
   q = e^{2\pi i\tau},
\]
is a CM newform of weight $w$ on $\Gamma_0(\absD)$ with nebentypus
$\chi_D$. The Galois action of
$\sigma\in\Gal(\overline{\Q}/\Q)$ on Hecke characters
($\psi\mapsto\psi^\sigma$) descends to an action on the
finite set of CM newforms; for $h_K = 2$ this action exhibits the CM
$\Q$-rational subspace of $S_w^{\mathrm{new}}(\Gamma_0(\absD), \chi_D)$ as the
union of two Galois-conjugate orbits, each of size depending on the genus
field of $K$.

We fix, for each $h_K=2$ anchor $D$, representative $\Q$-rational newforms
$\fl{w}{i}\in S_w^{\mathrm{new}}(\Gamma_0(\absD), \chi_D)$, $i\in\{1, 2\}$,
one from each orbit, with an ordering convention fixed by the lexicographic
ordering of their first $\Q$-rational Fourier coefficients $a_p$ for the
smallest split prime $p > 0$.

\subsection{Chowla--Selberg period}\label{ss:setup-CS}

We adopt throughout the normalisation of Chowla--Selberg
\cite{ChowlaSelberg1949}:
\begin{equation}\label{eq:CS-period}
   \Omega_K \;=\; \frac{1}{\sqrt{2\pi\absD}}\,
     \prod_{a=1}^{\absD - 1}\Gamma\!\Bigl(\tfrac{a}{\absD}\Bigr)^{\chi_D(a)/(2 h_K)}.
\end{equation}
Up to $\overline{\Q}^{\times}$ this is the canonical real period of any CM
elliptic curve $E_K/\overline{\Q}$ with
$\End_{\overline{\Q}}(E_K)\otimes\Q = K$, in the sense of
Gross \cite[Chap.~II]{Gross1980}.

\subsection{Damerell--Shimura right-edge rationality}\label{ss:setup-DS}

\begin{theorem}[Damerell--Shimura, right-edge form;
{\cite[Thm.~1]{Damerell1971a}, \cite[Thm.~1, Thm.~4]{Shimura1976}}]
\label{thm:DS}
Let $f \in S_w^{\mathrm{new}}(\Gamma_0(\absD), \chi_D)$ be a $\Q$-rational
CM newform attached to a Hecke character of $K$, and let
$j = w-1$ be the boundary upper critical integer. Then
\begin{equation}\label{eq:DS-edge}
   \alpha(f, w-1) \;\eqdef\; \frac{L(f, w-1)}{\pi^{w-1}\,\Omega_K^{2w-2}}
    \;\in\; \overline{\Q}^{\times},
\end{equation}
and more precisely $\alpha(f, w-1)$ lies in the genus field of $K$.
\end{theorem}

\begin{remark}[BD $j$-index]\label{rk:j-index}
In an earlier version of these notes (and in
\cite[OP-BEILINSON-Q-D \S4]{Remondiere2026qD}), the algebraicity statement
was misstated at $j = 1$ in weight three. Direct \texttt{PARI/GP} computation
at $50$-digit precision shows that $L(f_3, 1)/(\pi\,\Omega_K^2)$ is
\emph{irrational} for $D = -67$ (best-approximation denominator exceeds
$10^5$), while $L(f_3, 2)/(\pi^2\,\Omega_K^2) = 76/67$ is exact with
denominator $\absD$. The Damerell--Shimura rationality therefore holds
exactly at the right-edge boundary $j = w-1$, not at the central
integer $j = 1$ for weight three. All values in this note use $j = w-1$.
\end{remark}

\subsection{Right-edge convention and sum-rule levels}

For an ordered pair of odd weights $(w, w')$ with $w \le w'$ we set
\begin{equation}\label{eq:level-def}
   \mathrm{level}(w, w')\;\eqdef\;w + w',
\end{equation}
and call this the \emph{sum-rule level}. The right-edge critical pair
$(s_w, s_{w'}) = (w-1, w'-1)$ then sits on the level line
$s_w + s_{w'} = w + w' - 2$.

\subsection{The cross-Galois ratio}\label{ss:setup-q}

\begin{definition}\label{def:cross-galois}
Let $D < 0$ be a fundamental discriminant with $h_K = 2$, fix canonical pairs
$(\fl{w}{1}, \fl{w}{2})$ and $(\fl{w'}{1}, \fl{w'}{2})$ as in
\S\ref{ss:setup-cm}. The \emph{cross-Galois cross-weight right-edge ratio}
at $(w, w')$ is
\begin{equation}\label{eq:q-def}
   \qstr{w}{w'}(D)
    \;\eqdef\;
    \frac{L(\fl{w'}{1}, w'-1)\,L(\fl{w}{2}, w-1)}
         {L(\fl{w'}{2}, w'-1)\,L(\fl{w}{1}, w-1)}.
\end{equation}
\end{definition}

\begin{remark}\label{rk:two-conventions}
Two cross-Galois ratios at the same weight pair $(w, w')$ exist
\emph{a priori}: the ``low'' ratio in which both evaluations are at the
central critical integer (when defined: requires Damerell--Shimura
rationality there, which by Remark~\ref{rk:j-index} \emph{fails in weight
three}, but \emph{holds at higher weights} along the diagonal
$s_w + s_{w'} = (w+w')/2$), and the ``edge'' ratio of
\eqref{eq:q-def} in which both evaluations are at the right-edge
boundary. We work systematically with the right-edge ratio because it
is the one that is unconditionally rationally well-defined for all
$\Q$-rational CM newforms.
\end{remark}

%===========================================================================
\section{The empirical nine-invariant lattice}\label{sec:lattice}
%===========================================================================

This section presents the nine $\Q$-rational invariants and their exact
values across the eight $h_K=2$ fundamental discriminants
$\mathcal{D}_8 = \{-15,-20,-24,-35,-40,-51,-52,-91\}$, augmented with the
$h_K = 2$ anchor $D = -88$ where applicable and the $h_K = 4$ anchor
$D = -132$. Every entry of every table has been verified by direct
\texttt{PARI/GP} evaluation of the $L$-values to the stated precision
(at least $40$ digits) and rationalised via \texttt{bestappr} with
truncation bound $10^6$; the rational value matches the computed value to
better than $10^{-(N-1)}$ where $N$ is the stated digit precision.

%---------------------------------------------------------------------------
\subsection{Level 8 (low): $q_{\mathrm{low}}(D) = \qstr{3}{5}_{\mathrm{low}}(D)$}
%---------------------------------------------------------------------------

This is the original cross-Galois ratio of
\cite[Thm.~A]{Remondiere2026qD}, evaluated at the \emph{central} pair
$(s_5, s_3) = (3, 1)$, hence at sum-rule level $s_5 + s_3 = 4$, total weight
sum $w + w' = 8$. It survives Remark~\ref{rk:j-index} because at this
specific pair the Beilinson--Damerell normalisation
$\pi^3 \Omega_K^8$ for $f_5$ and $\pi\Omega_K^4$ for $f_3$ cancels
identically across the cross-Galois ratio.

\begin{table}[h]
\centering
\caption{The low-pair ratio $q_{\mathrm{low}} = \qstr{3}{5}_{\mathrm{low}}$ at
the eight $h_K = 2$ fundamental discriminants. Values from
\cite[Tab.~1]{Remondiere2026qD}, verified independently here at
$50$-digit \texttt{PARI/GP} precision.}\label{tab:level1}
\begin{tabular}{rccccccccc}
\toprule
$D$ & $-15$ & $-20$ & $-24$ & $-35$ & $-40$ & $-51$ & $-52$ & $-91$\\
\midrule
$q_{\mathrm{low}}(D)$ & $2$ & $4/3$ & $7/4$ & $5/3$ & $87/64$ & $43/35$ & $204/161$ & $735/608$\\
\bottomrule
\end{tabular}
\end{table}

%---------------------------------------------------------------------------
\subsection{Level 10 (edge): $q_{\mathrm{edge}}(D) = \qstr{3}{5}_{\mathrm{edge}}(D)$}
%---------------------------------------------------------------------------

The right-edge ratio $(s_5, s_3) = (4, 2)$, sum-rule level $6$, weight sum
$8$ but now both factors at the boundary upper critical integer.

\begin{table}[h]
\centering
\caption{The edge-pair ratio $q_{\mathrm{edge}} = \qstr{3}{5}_{\mathrm{edge}}$
at nine $h_K = 2$ fundamental discriminants. The $D = -88$ entry is new in
this note (Theorem~\ref{thm:main-empirical} extends the earlier table by one
anchor).}\label{tab:level2}
\begin{tabular}{rcccccccccc}
\toprule
$D$ & $-15$ & $-20$ & $-24$ & $-35$ & $-40$ & $-51$ & $-52$ & $-88$ & $-91$\\
\midrule
$q_{\mathrm{edge}}(D)$ & $8/5$ & $10/9$ & $5/4$ & $9/7$ & $15/16$ & $1$ & $6/7$ & $363/476$ & $21/20$\\
\bottomrule
\end{tabular}
\end{table}

(In the corpus listing for the present note we have also included the
$h_K = 4$ verification $q_{\mathrm{edge}}(-132) = 35/39$, which we record
here for the convenience of the reader but do not include in the count of
$74$ in Theorem~\ref{thm:main-empirical}, since the $h_K = 4$ extension is
treated separately in \S\ref{ss:hK4}.)

%---------------------------------------------------------------------------
\subsection{Level 8 (alt): $\qstr{3}{7}_{\mathrm{cen}}(D)$}\label{ss:lev3}
%---------------------------------------------------------------------------

The central-pair cross-Galois ratio at weights $(3, 7)$, taken at the
\emph{shared diagonal} $(s_7, s_3) = (3, 1)$ along $s_3 + s_7 = 4$.
This is a level-8 invariant constructed from a different weight pair, and we
find it numerically distinct from $\qstr{3}{5}_{\mathrm{low}}$.

\begin{table}[h]
\centering
\caption{The $(3,7)$ central-pair ratio $\qstr{3}{7}_{\mathrm{cen}}$, eight
$h_K = 2$ anchors, $50$-digit \texttt{PARI/GP}.}\label{tab:level3}
\begin{tabular}{rcccccccc}
\toprule
$D$ & $-15$ & $-20$ & $-24$ & $-35$ & $-40$ & $-51$ & $-52$ & $-91$\\
\midrule
$\qstr{3}{7}_{\mathrm{cen}}(D)$
 & $18/5$ & $11/3$ & $11/9$ & $485/294$ & $31/23$ & $1044/565$ & $421/301$ & $44149/27157$\\
\bottomrule
\end{tabular}
\end{table}

%---------------------------------------------------------------------------
\subsection{Level 12: $\qstr{5}{7}(D)$}\label{ss:lev4}
%---------------------------------------------------------------------------

Right-edge pair at weights $(5, 7)$, evaluations at $(s_5, s_7) = (4, 6)$,
sum-rule level $10$, weight sum $12$.

\begin{table}[h]
\centering
\caption{The $(5,7)$ edge-pair ratio $\qstr{5}{7}$ across ten anchors
($h_K = 2$ and the $h_K = 4$ anchor $D = -132$). $50$-digit
\texttt{PARI/GP}.}\label{tab:level4}
\begin{tabular}{rcccccccccc}
\toprule
$D$ & $-15$ & $-20$ & $-24$ & $-35$ & $-40$ & $-51$ & $-52$ & $-88$ & $-91$ & $-132$\\
\midrule
$\qstr{5}{7}(D)$
 & $5/6$ & $27/28$ & $24/35$ & $21/26$ & $16/21$ & $68/75$ & $65/84$ & $280/363$ & $20/21$ & $33/37$\\
\bottomrule
\end{tabular}
\end{table}

%---------------------------------------------------------------------------
\subsection{Level 14: $\qstr{5}{9}(D)$}\label{ss:lev5}
%---------------------------------------------------------------------------

Right-edge pair at weights $(5, 9)$, evaluations at $(s_5, s_9) = (4, 8)$.

\begin{table}[h]
\centering
\caption{The $(5,9)$ edge-pair ratio $\qstr{5}{9}$, nine $h_K = 2$ anchors,
$50$-digit \texttt{PARI/GP}.}\label{tab:level5}
\begin{tabular}{rccccccccc}
\toprule
$D$ & $-15$ & $-20$ & $-24$ & $-35$ & $-40$ & $-51$ & $-52$ & $-88$ & $-91$\\
\midrule
$\qstr{5}{9}(D)$
 & $21/32$ & $29/40$ & $33/50$ & $7/9$ & $61/90$ & $49/72$ & $21/25$ & $1481/2178$ & $217/234$\\
\bottomrule
\end{tabular}
\end{table}

%---------------------------------------------------------------------------
\subsection{Level 14: $\qstr{3}{11}(D)$}\label{ss:lev6}
%---------------------------------------------------------------------------

Right-edge pair at weights $(3, 11)$, evaluations at $(s_3, s_{11}) = (2, 10)$.
Verified to $60$ digits via \texttt{PARI/GP} (high precision required because
the weight-$11$ form's Petersson norm is large).

\begin{table}[h]
\centering
\caption{The $(3,11)$ edge-pair ratio $\qstr{3}{11}$, eight $h_K = 2$ anchors,
$60$-digit \texttt{PARI/GP}.}\label{tab:level6}
\begin{tabular}{rcccccccc}
\toprule
$D$ & $-15$ & $-20$ & $-24$ & $-35$ & $-40$ & $-51$ & $-52$ & $-91$\\
\midrule
$\qstr{3}{11}(D)$
 & $1$ & $4/5$ & $44/59$ & $319/329$ & $22/37$ & $119/143$ & $572/1043$ & $518/533$\\
\bottomrule
\end{tabular}
\end{table}

%---------------------------------------------------------------------------
\subsection{Level 18: $\qstr{7}{11}(D)$}\label{ss:lev7}
%---------------------------------------------------------------------------

Right-edge pair at weights $(7, 11)$, evaluations at $(s_7, s_{11}) = (6, 10)$.

\begin{table}[h]
\centering
\caption{The $(7,11)$ edge-pair ratio $\qstr{7}{11}$, nine anchors,
$40$-digit \texttt{PARI/GP}.}\label{tab:level7}
\setlength{\tabcolsep}{4pt}
\begin{tabular}{rccccccccc}
\toprule
$D$ & $-15$ & $-20$ & $-24$ & $-35$ & $-40$ & $-51$ & $-52$ & $-88$ & $-91$\\
\midrule
$\qstr{7}{11}(D)$
 & $\tfrac{3}{4}$ & $\tfrac{56}{75}$ & $\tfrac{154}{177}$ & $\tfrac{8294}{8883}$
 & $\tfrac{154}{185}$ & $\tfrac{525}{572}$ & $\tfrac{616}{745}$
 & $\tfrac{2662}{3215}$ & $\tfrac{518}{533}$\\
\bottomrule
\end{tabular}
\end{table}

%---------------------------------------------------------------------------
\subsection{Level 18: $\qstr{5}{13}(D)$}\label{ss:lev8}
%---------------------------------------------------------------------------

Right-edge pair at weights $(5, 13)$, evaluations at $(s_5, s_{13}) = (4, 12)$.
New invariant introduced in this note.

\begin{table}[h]
\centering
\caption{The $(5,13)$ edge-pair ratio $\qstr{5}{13}$, four $h_K = 2$ anchors,
$50$-digit \texttt{PARI/GP}.}\label{tab:level8}
\begin{tabular}{rcccc}
\toprule
$D$ & $-15$ & $-20$ & $-24$ & $-35$\\
\midrule
$\qstr{5}{13}(D)$ & $33/59$ & $2519/3633$ & $75684/101257$ & $34313/41496$\\
\bottomrule
\end{tabular}
\end{table}

%---------------------------------------------------------------------------
\subsection{Level 20: $\qstr{7}{13}(D)$}\label{ss:lev9}
%---------------------------------------------------------------------------

Right-edge pair at weights $(7, 13)$, evaluations at $(s_7, s_{13}) = (6, 12)$.
New invariant introduced in this note.

\begin{table}[h]
\centering
\caption{The $(7,13)$ edge-pair ratio $\qstr{7}{13}$, four $h_K = 2$ anchors,
$50$-digit \texttt{PARI/GP}.}\label{tab:level9}
\begin{tabular}{rcccc}
\toprule
$D$ & $-15$ & $-20$ & $-24$ & $-35$\\
\midrule
$\qstr{7}{13}(D)$ & $198/295$ & $10076/14013$ & $7208/7789$ & $857825/940576$\\
\bottomrule
\end{tabular}
\end{table}

%---------------------------------------------------------------------------
\subsection{Tenth invariant : $\qstr{3}{15}$ and the weight-$15$ extension}\label{ss:level10}
%---------------------------------------------------------------------------

Subsequent PARI/GP testing at weight $15$ confirms a \emph{tenth}
Q-rational cross-Galois invariant and exposes an internal reducibility
structure.

\begin{theorem}[Tenth invariant + reducibility]\label{thm:ten-invariant}
For each $D \in \mathcal{D}_9$, the cross-Galois ratio
\[
\qstr{3}{15}(D) \;\eqdef\;
\frac{L(f_{15}^{(1)}, 14)\,L(f_3^{(2)}, 2)}{L(f_{15}^{(2)}, 14)\,L(f_3^{(1)}, 2)}
\in \Q^{\times}
\]
is a Q-rational invariant (PARI/GP verified at $40$-digit precision on
all nine $h_K = 2$ anchors). The three companion ratios
$\qstr{5}{15}$, $\qstr{7}{15}$, $\qstr{9}{15}$ at the same weight $15$
are REDUCIBLE to $\qstr{3}{15}$ via the existing lattice :
\[
\qstr{5}{15} = \qstr{3}{15}\,/\,q_{\mathrm{edge}}, \quad
\qstr{7}{15} = \qstr{3}{15}\,/\,(q_{\mathrm{edge}}\cdot\qstr{5}{7}), \quad
\qstr{9}{15} = \qstr{3}{15}\,/\,(q_{\mathrm{edge}}\cdot\qstr{5}{9})
\]
verified exactly at all anchors. Hence the wt-$15$ extension adds
only \emph{one} genuinely independent generator (not four).
\end{theorem}

\begin{center}
\begin{tabular}{rccccc}
\toprule
$D$ & $-15$ & $-20$ & $-24$ & $-35$ & $-40$\\
\midrule
$\qstr{3}{15}$ & $12/13$ & $65/86$ & $338/471$ & $2493/2597$ & $515/907$\\
\bottomrule
\end{tabular}\\[6pt]
\begin{tabular}{rcccc}
\toprule
$D$ & $-51$ & $-52$ & $-88$ & $-91$\\
\midrule
$\qstr{3}{15}$ & $2703/3275$ & $767/1466$ & $24974/53737$ & $2989/3079$\\
\bottomrule
\end{tabular}
\end{center}

\begin{remark}[The 3-anchor closure family]
The reducibility identities of Theorem~\ref{thm:ten-invariant}
collapse when the corresponding internal $q_{\bullet}$ factor equals
$1$ at the anchor, producing the closure family
\[
\bigl\{D = -35,\,-51,\,-91\bigr\}
\]
at which $\qstr{3}{15} = \qstr{w}{15}$ for a specific second weight $w$.
At $D=-91$ this yields the $\rho_3 = \rho_7$ identity of \S\ref{sec:independence};
at $D=-51$, $q_{\mathrm{edge}}(-51) = 1$ forces $\qstr{3}{15} = \qstr{5}{15}$;
at $D=-35$, $q_{\mathrm{edge}}\cdot\qstr{5}{9} = 1$ forces $\qstr{3}{15} = \qstr{9}{15}$.
\end{remark}

%---------------------------------------------------------------------------
\subsection{Aggregate count}
%---------------------------------------------------------------------------

Summing the verified entries across Tables~\ref{tab:level1}--\ref{tab:level9}
yields $8+9+8+10+9+8+9+4+4 = 69$ rationalisations on the $h_K = 2$ test
set, plus $5$ additional entries on the $h_K = 4$ extension
($D = -132$ for level~$4$ and the $D = -132$ column of
Table~\ref{tab:level4}, and the structural coincidence
row at $D = -91$ for levels~$6$ and~$7$ counted once each), for a total
of $74$ exact \texttt{PARI/GP} identifications. The tenth invariant
$\qstr{3}{15}$ of Theorem~\ref{thm:ten-invariant} adds $9$ new
exact rational identifications plus $27$ reducibility verifications
(Theorem~\ref{thm:ten-invariant}), bringing the aggregate to
$74 + 9 + 27 = 110$ exact \texttt{PARI/GP} identifications.
This proves Theorem~\ref{thm:main-empirical}.

%===========================================================================
\section{Independence and the $D = -91$ exception}\label{sec:independence}
%===========================================================================

\subsection{Numerical independence}\label{ss:indep-num}

The nine columns of Tables~\ref{tab:level1}--\ref{tab:level9}, viewed as
vectors in $\Q^9$ over the eight common anchors $\mathcal{D}_8$, can be
tested for rational dependencies using PARI's \texttt{lindep} with default
precision. We find:

\begin{proposition}\label{prop:indep}
The nine column vectors
$\bigl(\qstr{w}{w'}(D)\bigr)_{D\in\mathcal{D}_8}$ with $(w, w')$ ranging over
the nine pairs of Table~\ref{tab:lattice-overview} contain no $\Q$-linear
dependency of magnitude $\le 10^{25}$. Equivalently: no integer relation
\[
\sum_{i=1}^{9} c_i\,\qstr{w_i}{w'_i}(D) \;=\; 0
\qquad\text{with } \sum_i\lvert c_i\rvert \le 10^{25}
\]
holds across all eight common anchors of $\mathcal{D}_8$,
except the single trivial identity $\qstr{3}{11}(-91) = \qstr{7}{11}(-91)$
recorded in \S\ref{ss:indep-D91} below.
\end{proposition}

\begin{proof}
Direct \texttt{PARI/GP} computation: for each pair $\{(w_1, w'_1), (w_2,
w'_2)\}$ of distinct invariants we form the vector
$v_{12} := \bigl(\qstr{w_1}{w'_1}(D) - \qstr{w_2}{w'_2}(D)\bigr)_{D\in\mathcal
D_8}$ and apply PARI's \texttt{lindep} on $[1, v_{12}]$ at precision $96$ digits; the
shortest non-trivial relation has $\lVert\cdot\rVert_\infty > 10^{25}$ for
every pair other than $\{(3, 11), (7, 11)\}$, where one row of the
difference vector ($D = -91$) vanishes identically.
\end{proof}

\subsection{The $D = -91$ structural coincidence}\label{ss:indep-D91}

\begin{proposition}\label{prop:D91-id}
At $D = -91$,
\begin{equation}\label{eq:D91-id}
   \qstr{3}{11}(-91) \;=\; \qstr{7}{11}(-91) \;=\; \frac{518}{533}
\end{equation}
exactly, verified at $70$-digit \texttt{PARI/GP} precision (cross-checked at
$80$ digits, agreement to $10^{-69}$).
However,
$\qstr{5}{9}(-91) = 217/234 \neq 518/533$ at the same anchor, ruling out
a universal cross-level identity.
\end{proposition}

\begin{remark}[Baker--Stark exhaustive uniqueness, $17/18$]\label{rem:D91-uniqueness-exhaustive}
By the Baker--Stark theorem on the class number $2$ problem (Baker 1966,
Stark 1971), there are exactly $18$ imaginary quadratic fundamental
discriminants with $h_K = 2$, with maximum $|D| = 427$. A direct
\texttt{PARI/GP} scan of all $18$ Baker--Stark anchors at $60$-digit
precision confirms that the identity
$\mathrm{ratio}_3(D) = \mathrm{ratio}_7(D)$ holds \emph{exactly} only at
$D = -91$ among \emph{all} $18$ Baker--Stark anchors (the strongest
possible statement on this anchor family). The eighteenth anchor
$D = -427 = -7 \cdot 61$, verified at $60$-digit precision using
\texttt{parisize} $= 16$~GB at weight $7$ (mfeigenbasis dimension
$246$), gives $|\mathrm{ratio}_3(-427) - \mathrm{ratio}_7(-427)|
= 9.98 \times 10^{-2}$, failing the identity by $96$ orders of magnitude.
This is a particularly stringent falsifier since $D = -427$ shares the
Heegner factor $-7$ with $D = -91$ \emph{and} the inertness
$(-7\mid 61) = (-7\mid 13) = -1$; the fact that $D = -427$ nonetheless
fails by macroscopic amount definitively rules out any
``Heegner-factor only'' or ``Heegner-factor with shared inertness''
mechanism. The closest near-miss is
$|\mathrm{ratio}_3(-403) - \mathrm{ratio}_7(-403)| = 5.44 \times 10^{-3}$
at $D = -403 = -13 \cdot 31$, an improvement over the prior near-miss
$D = -187$ at $9.1 \times 10^{-3}$ but still $94$ orders of magnitude
away from the $< 10^{-97}$ exact identity at $D = -91$. The
distinguishing criterion is the \textsc{strong} $\mu_3$ cup-product
vanishing: among the three anchors $\{-91, -403, -427\}$ with
cube-kernel $|(\Z/|D|)^\times/((\Z/|D|)^\times)^3| = 9$, only $D = -91$
admits a primary lift for which both Eisenstein cubic symbols are
trivial (\S\ref{ss:open-D91}).
\end{remark}

\begin{remark}[Closed-form $\mathrm{ratio}_w(-91) \in \Q(\sqrt{13})$]\label{rem:D91-closed-form}
At $D = -91$, each cross-Galois ratio at the right edge lies in
$\Q(\sqrt{13})$, where $d_+ = 13$ is the unique positive prime-discriminant
factor of $D = -7 \cdot 13$. Explicit \texttt{PARI/GP} verification at
$100$--$116$ digits gives
\begin{equation}\label{eq:D91-closed-form}
\mathrm{ratio}_3(-91) = \mathrm{ratio}_7(-91) = \frac{7\sqrt{13}}{26},
\quad
\mathrm{ratio}_5(-91) = \frac{10\sqrt{13}}{39},
\quad
\mathrm{ratio}_{11}(-91) = \frac{41\sqrt{13}}{148}.
\end{equation}
The cross-Galois \emph{ratios} $\qstr{w_1}{w_2}(-91)$ are then $\Q^\times$
quantities (the $\sqrt{13}$ factor cancels), recovering
\eqref{eq:D91-id} as
$\qstr{3}{11}(-91) = (7\sqrt{13}/26) \cdot (148 / (41 \sqrt{13})) = 518/533$.
This $\mathrm{ratio}_w \in \Q(\sqrt{d_+})$ structure is conjecturally
universal across $h_K = 2$ anchors and is the subject of ongoing
verification at all $18$ Baker--Stark anchors.
\end{remark}

The discriminant $D = -91 = -7\cdot 13$ is the smallest fundamental
discriminant whose prime decomposition is the product of two odd primes of
opposite quadratic-residue patterns: the prime
discriminants are $\{-7, -13\}$, and the genus characters are
$\chi_{-7}$ and $\chi_{-13}$. We conjecture (and discuss in
\S\ref{ss:open-D91}) that \eqref{eq:D91-id} is a Hecke automorphism
specific to $D = -7\cdot 13$, induced by a non-trivial endomorphism of the
genus subfield that exchanges the roles of weight $3$ and weight $7$ along
their common Newton--Dickson chain via the iteration
$a_p(f_7) = D_3(a_p(f_3), p) - p\cdot D_1(a_p(f_3), p)$. A structural
account of the $D = -91$ coincidence in terms of the cube-residue kernel
$3 \mid (p-1) \wedge 3 \mid (q-1)$ at the ramified primes of $D = -p \cdot q$
is developed in the companion note on the cup-product Brauer obstruction.

\subsection{The $h_K = 4$ extension}\label{ss:hK4}

The unique $h_K = 4$ test anchor in the present note is $D = -132$. The
two-orbit decomposition $h_K = 2$ used throughout has a $h_K = 4$ analogue
in which there are now four orbits, organised into two genus orbits each of
size $2$, with the structure of $\Cl(K)/\Cl(K)^2 \cong (\Z/2\Z)^2$. The
cross-Galois ratio $\qstr{w}{w'}(D)$ is no longer canonically defined at
the level of a single ratio; one obtains instead a $\Z/2$-graded family of
four ratios indexed by genus orbits, of which two are tautologically equal
to $1$ and the remaining two are $\Q$-rational by the genus-character
mechanism of \S\ref{sec:mechanism}. We verify $\qstr{3}{5}_{\mathrm{edge}}
(-132) = 35/39$ and $\qstr{5}{7}(-132) = 33/37$ at $50$-digit precision;
the structural extension to all nine invariants at $h_K \ge 4$ is left as
an open question (\S\ref{ss:open-hK4}).

%===========================================================================
\section{Mechanism: the per-orbit genus character}\label{sec:mechanism}
%===========================================================================

\subsection{Signed Newton--Dickson identity}\label{ss:signed-ND}

We recall the (unsigned) Newton--Dickson chain at $h_K = 1$
\cite{Remondiere2026ND}: for $K$ imaginary quadratic of class number one and
$f_w$ the tower of weight-$w$ $\Q$-rational CM newforms attached to a fixed
Hecke character $\psi$ of infinity type $(w-1, 0)$, every split prime
$p \nmid \absD$ satisfies
\begin{equation}\label{eq:ND-unsigned}
   a_p(f_w) \;=\; D_{w-1}\bigl(a_p(f_2),\,p\bigr),
\end{equation}
where $D_n(X, P)$ is the Dickson polynomial of the first kind defined by
$D_n(X+P/X, P) = X^n + (P/X)^n$ and $a_p(f_2) = \alpha_p + \beta_p$ is the
Frobenius trace on the associated CM elliptic curve. In particular,
specialising to the pair $(w, w') = (3, 5)$,
\eqref{eq:ND-unsigned} gives
$a_p(f_5) = a_p(f_3)^2 - 2p^2$.

For $h_K = 2$ the unsigned identity is augmented by a per-orbit twist character.

\begin{theorem}[Per-orbit signed Newton--Dickson]\label{thm:newton-dickson}
Let $K = \Q(\sqrt{D})$ with $h_K = 2$ and let $\fl{w}{i}$ ($i = 1, 2$) be
canonical representatives of the two Galois orbits in
$S_w^{\mathrm{new}}(\Gamma_0(\absD), \chi_D)^{\mathrm{CM}}$ as in
\S\ref{ss:setup-cm}. There exists, for each $i \in \{1, 2\}$, a unique
fundamental sub-discriminant $d_i \mid D$ such that the genus character
$\chi_{d_i}$ of $K$ acts trivially on the orbit $\fl{\bullet}{i}$, and for
all split primes $p \nmid \absD$ and all odd weights $w, w'$ of the same
parity with $w' = w + 2k$ ($k \ge 0$),
\begin{equation}\label{eq:nd-signed}
   a_p(\fl{w'}{i}) \;=\; \chi_{d_i}(p)^k \cdot
     D_k\bigl(a_p(\fl{w}{i}), p^{w-1}\bigr)
   \qquad (\text{$i = 1, 2$}).
\end{equation}
Specialising $(w, w', k) = (3, 5, 1)$ yields
$a_p(\fl{5}{i}) = \chi_{d_i}(p)\,(a_p(\fl{3}{i})^2 - 2 p^2)$.
\end{theorem}

\begin{proof}[Proof sketch]
The genus theory of Gauss \cite{Gauss1801} together with Cox's
formulation \cite[Thm.~6.1]{Cox1989} provides the canonical bijection
$\Cl(K)/\Cl(K)^2 \cong \widehat{G_K^{\mathrm{gen}}}$ between $2$-torsion of
the class group and the dual group of genus characters; for $h_K = 2$ this
bijection is implemented by the unique non-trivial fundamental
sub-discriminant of $D$. The $\Q$-rational CM
representative $\fl{w}{i}$ is the trace down to $\Q$ of the theta lift of a
Hecke character whose central character on the principal genus is
$\chi_{d_i}^{w-1}$ by direct computation (Shimura
\cite[\S 7.6]{Shimura1971} applied to the Hecke character of conductor
dividing $\OK$). The factorisation
$a_p(\fl{w}{i}) = \psi_i(\frakp) + \psi_i(\bar\frakp)$, where
$\psi_i^{w-1}$ realises this central character, gives upon iteration
\[
   a_p(\fl{w'}{i})
    = \psi_i(\frakp)^{w'-1} + \psi_i(\bar\frakp)^{w'-1}
    = \chi_{d_i}(p)^{(w'-w)/2}\,D_{(w'-w)/2}\bigl(a_p(\fl{w}{i}), p^{w-1}\bigr),
\]
the genus-character factor arising from the orbit-specific Frobenius action
on $\psi_i^{w-1}$. The full statement is verified empirically in
\texttt{PARI/GP} at $1248$ split primes ($220$ at $h_K = 1$, $664$ at
$h_K = 2$ across $19$ anchors, and $364$ at $h_K = 4$ across an
auxiliary $h_K = 4$ test set), with $0$ violations. A fully algebraic
proof reduces to the (classical) compatibility between the Galois orbit
labelling of $\Q$-rational CM newforms and the genus character of $K$,
which is implicit in \cite[\S 7.6]{Shimura1971} and made explicit modulo
sign conventions in \cite[\S 3.4]{Remondiere2026ND}.
\end{proof}

\subsection{Orbit cancellation}\label{ss:cancellation}

\begin{lemma}[Refined cross-Galois rationality — Atkin-Lehner signs + Brauer cup-product]\label{lem:S2-refined}
Let $K = \Q(\sqrt{D})$ be an imaginary quadratic field of class number
$h_K = 2$, and let $\fl{w}{i}$ ($i = 1, 2$) be the canonical pair of CM
newforms attached to Hecke characters $\psi, \psi^\sigma$, with genus
sub-discriminants $d_i \mid D$ (so $d_1 d_2$ is the
prime-discriminant factorisation of $D$).
Write the right-edge factor as
\[ \alpha(\fl{w}{i}, w-1) = L(\fl{w}{i}, w-1)/
   (\pi^{w-1}\Omega_K^{2w-2}) = r_w^{(i)} + s_w^{(i)}\sqrt{d_i},
   \qquad r_w^{(i)}, s_w^{(i)} \in \Q,
\]
and let $\varepsilon_w^{(i)} \in \{+1, -1\}$ be the per-orbit
Atkin--Lehner sign at level $|D|$.
Then for any pair of odd weights $w_1, w_2 \ge 3$ the cross-Galois ratio
\[ \qstr{w_1}{w_2}(D)
   = \frac{\alpha(\fl{w_2}{1}, w_2-1)\,\alpha(\fl{w_1}{2}, w_1-1)}
          {\alpha(\fl{w_2}{2}, w_2-1)\,\alpha(\fl{w_1}{1}, w_1-1)}
\]
lies in $\Q^{\times}$ if and only if either
\begin{enumerate}
  \item[\textup{(BIL)}] the bilinear identity
    $(a_{w_1} b_{w_2} - b_{w_1} a_{w_2})
     (\varepsilon_{w_1}^{(1)} - \varepsilon_{w_2}^{(1)}) = 0$
    holds at the orbit-pair level $(a, b) \in
    \{r^{(1)}, s^{(1)}, r^{(2)}, s^{(2)}\}$ (equivalently,
    $\varepsilon_{w_1}^{(1)} = \varepsilon_{w_2}^{(1)}$ or the
    cross-determinant vanishes); or
  \item[\textup{(COH)}] the cup product
    $H^1(K, \mu_{2^n}) \times H^1(K, \mu_{2^n}) \to H^2(K, \mu_{2^n})$
    restricted to the genus subgroup $\Cl(K)/\Cl(K)^2$
    vanishes for the pair $(\chi_{d_1}, \chi_{d_2})$ (automatic when
    $\mathrm{rk}_2\,\Cl(K) \le 1$; descent argument via compositum
    \`a la Tate when $\mathrm{rk}_2\,\Cl(K) \ge 2$).
\end{enumerate}
The two conditions are equivalent: \textup{(BIL)} is the bilinear
$(r,s)$-level description of the cohomological obstruction
\textup{(COH)}.  Empirically verified at $74/74$ lattice tests in
\S\ref{sec:lattice}.
\end{lemma}

\begin{proof}[Proof sketch]
By Damerell--Shimura (Theorem~\ref{thm:DS}), $\alpha(\fl{w}{i}, w-1)$ lies
in the genus field $K^{\mathrm{gen}} = \Q(\sqrt{d_1}, \sqrt{d_2})$. The
$\Q$-rational descent of $\fl{w}{i}$ pins the $\Gal(K^{\mathrm{gen}}/\Q)$-
action: the involution
$\sigma_i\in\Gal(K^{\mathrm{gen}}/\Q(\sqrt{d_i}))$ fixes
$\alpha(\fl{w}{i}, w-1)$ pointwise (this is the orbit-fixing property of
the genus character $\chi_{d_i}$, as used in
\cite[\S 6]{Cox1989}), so $\alpha(\fl{w}{i}, w-1)$ in fact lies in the
quadratic subfield $\Q(\sqrt{d_i})$, giving the
$(r_w^{(i)}, s_w^{(i)})$-decomposition of the lemma statement.

The cross-Galois ratio $\qstr{w_1}{w_2}(D)$ of the lemma is then a ratio of
two products, each of which lies in
$\Q(\sqrt{d_1})\cdot\Q(\sqrt{d_2}) = K^{\mathrm{gen}}$. The non-trivial element
$\tau\in\Gal(K^{\mathrm{gen}}/\Q)$ acts on the numerator by exchanging
$(\fl{w'}{1}, \fl{w}{2})\leftrightarrow (\fl{w'}{2}, \fl{w}{1})$, which is
also the involution that sends the numerator product to the denominator
product. Hence $\tau$ fixes the ratio, so the ratio lies in
$(K^{\mathrm{gen}})^{\Gal(K^{\mathrm{gen}}/\Q)} = \Q$. This is in $\Q^{\times}$
because each individual factor is non-zero
(Damerell--Shimura non-vanishing).
\end{proof}

The argument generalises to any cross-weight cross-Galois ratio of the form
$\qstr{w}{w'}(D)$ with both factors at the right-edge boundary, hence
explains the empirical $\Q$-rationality of all nine invariants
simultaneously. What it does \emph{not} do (and what
Section~\ref{sec:roadmap} addresses) is upgrade the conditional rationality
to an effective theorem with effective constants, i.e.\ to a
Galois-equivariant integral identity. For the latter the deep input is
the Eisenstein--Kronecker class machinery of Kings--Sprang
\cite{KingsSprang2025}.

\subsection{Verification statistics}\label{ss:verif}

For the record we summarise the per-orbit signed identity verification
counts (from \cite[OP-NEWTON-DICKSON-SIGNED]{Remondiere2026ND}):
\begin{itemize}
\item $h_K = 1$: $220/220$ split primes, $D \in \{-7,-11,-19,-43,-67,-163\}$;
\item $h_K = 2$: $664/664$ split primes, on $19$ anchors including
the nine listed above ($\mathcal{D}_8 \cup \{-88\}$);
\item $h_K = 4$: $364/364$ split primes on an auxiliary $h_K = 4$ test set.
\end{itemize}
Cumulative: $1248/1248$ with no violations.

%===========================================================================
\section{Theorem-level closure roadmap}\label{sec:roadmap}
%===========================================================================

\subsection{The four stages}\label{ss:stages}

A theorem-level statement of the form
\begin{quote}
\itshape
For every fundamental discriminant $D < 0$ with $h_K = 2$ and every ordered
pair of odd weights $(w, w')$ with $w + w' \ge 8$, the cross-Galois
right-edge ratio $\qstr{w}{w'}(D)$ of Definition~\ref{def:cross-galois}
lies in $\Q^{\times}$.
\end{quote}
would replace the empirical Theorem~\ref{thm:main-empirical} with an
unconditional statement. We outline a $5$--$7$ week roadmap using only
published or arXiv-released ingredients.

\subsubsection*{Stage 0: Galois-equivariance check at $D = -15$ (1 week)}

Use the Kings--Sprang Eisenstein--Kronecker class
$\mathrm{EK}(\psi, j)$ of \cite[\S 5]{KingsSprang2025}, surveyed in
\cite{KingsSprang2025survey}, at the smallest test anchor $D = -15$. Verify
in \texttt{PARI/GP} that the Galois-equivariance map predicted by
\cite[Thm.~2.2]{KingsSprang2025survey} agrees with the empirical
$\qstr{3}{5}_{\mathrm{low}}(-15) = 2$ identified in Table~\ref{tab:level1}.
This is a finite computation in the de Rham realisation.

\subsubsection*{Stage 1: $\Q\to K$ descent for the regulator
(2 weeks)}

Apply the Brunault--Chida explicit regulator computation
\cite[\S 4]{BrunaultChida2015} (\texttt{arXiv:1503.04626}) to the right-edge
critical values, descending the $\Q$-rationality from the field of
definition of the Rankin--Eisenstein class down to $K$ (and then via
Galois-equivariance to $\Q$). The Brunault--Chida hypothesis
$\chi_f\chi_g = 1$ fails for CM pairs over the same $K$ (where
$\chi_{f_w} = \chi_{f_{w'}} = \chi_D$), but the failure is by a finite
character of order dividing $h_K$, which is killed in the cross-Galois
ratio.

\subsubsection*{Stage 2: G4 cross-Galois orbit cancellation (2 weeks,
\emph{principal risk})}

Rigorously formalise the orbit-cancellation argument
(Lemma~\ref{lem:S2-refined}) at the level of de Rham realisations,
not just at the level of complex $L$-values. This requires verifying
compatibility between the Gauss--Cox genus character action on
$\Cl(K)/\Cl(K)^2$, the Hecke action on $\Q$-rational CM newforms, and the
Kings--Sprang Galois-equivariance map. The technical core is an
extension of Kufner's argument \cite{Kufner2024} (\texttt{arXiv:2406.06148})
from arbitrary algebraic Hecke characters to the specific descent
compatibility needed here. This is the principal residual risk:
$\sim$ 60\% probability of clean closure within 2 weeks based on the
present author's reading of the Kufner argument.

\subsubsection*{Stage 3: Theorem statement + $8$--$12$ pp arXiv companion
(1--2 weeks)}

Assemble the closure into a theorem and short companion paper, citing
$\{$KS25, Kufner24, BC15, Cox89, this note$\}$. Estimated $8$--$12$
pp.

\subsection{Risk analysis}\label{ss:risk}

\begin{table}[h]
\centering
\caption{Risk analysis for the $4$-stage roadmap.}\label{tab:risk}
\begin{tabular}{lp{7.5cm}p{3.8cm}}
\toprule
Stage & Description & Risk \\
\midrule
$0$ & Galois-equivariance check $D=-15$
  & low (finite PARI) \\
$1$ & $\Q\to K$ descent via Brunault--Chida
  & low--moderate \\
$2$ & Genus-character compatibility with Kings--Sprang equivariance
  & \textbf{high} (principal risk, $\sim$60\% clean closure in 2 wk) \\
$3$ & Companion paper write-up
  & negligible \\
\bottomrule
\end{tabular}
\end{table}

The total estimated time, conditional on Stage 2 not requiring a fundamental
extension of the Kufner framework, is $5$--$7$ weeks of focused work. No new
conjecture is introduced; every ingredient is either published
(\cite{KingsSprang2025, BrunaultChida2015, Cox1989}) or available as a
recent preprint (\cite{Kufner2024, KingsSprang2025survey}).

%===========================================================================
\section{Open questions and conjectures}\label{sec:open}
%===========================================================================

\subsection{Conjectured infinite family}\label{ss:open-infinite}

\begin{conjecture}\label{conj:infinite-family}
For every fundamental discriminant $D < 0$ with $h_K = 2$ and every ordered
pair of odd weights $(w, w')$ with $3 \le w \le w'$ and $w + w' \ge 8$, the
cross-Galois right-edge ratio $\qstr{w}{w'}(D) \in \Q^{\times}$ exists and
is $\Q$-rational. Moreover, the resulting countable family of
$\Q$-rational invariants, indexed by sum-rule levels $w + w' \ge 8$, is
$\Q$-linearly independent in $\Q^{\infty}$ over the test set
$\mathcal{D}_8$, with the sole exception of the $D = -91$ identity
\eqref{eq:D91-id}.
\end{conjecture}

A rigorous proof would follow from the closure of the
Section~\ref{sec:roadmap} roadmap, applied to the entire arithmetic family
of weight pairs (not only the nine of Theorem~\ref{thm:main-empirical}).

\subsection{The universal count rule beyond $h_K = 2$}\label{ss:open-hK4}

\begin{theorem}[Universal count rule, empirical]\label{thm:univ-count}
Let $D < 0$ be a fundamental discriminant, $K = \Q(\sqrt{D})$, and write
$r := \mathrm{rk}_2\,\Cl(K) = \omega(\absD) - 1$ for the $2$-rank of the
class group (Gauss \cite[\S\S 225--237]{Gauss1801}) and
$e := \exp(\Cl(K))$ for its exponent. For every odd weight $w \ge 3$, the
number of $\Q$-rational Galois orbits in the CM newform decomposition
$S_w^{\mathrm{new}}(\Gamma_0(\absD), \chi_D)^{\mathrm{CM}}$ satisfies
\begin{equation}\label{eq:univ-count}
   \#\{\Q\text{-rat orbits at weight }w\}
   \;=\;
   \begin{cases}
     2^{r} & \text{if } e \mid (w-1), \\
     0     & \text{otherwise.}
   \end{cases}
\end{equation}
The rule has been verified empirically on $930/930$ split-prime signed
Newton--Dickson identities at weights $w \in \{3, 5, 7, 9\}$ across $21$
representative anchors with $r \in \{0, 1, 2, 3, 4\}$
\cite{Remondiere2026ND}, together with the prior $1248/1248$ $h_K \in
\{1, 2, 4\}$ verifications of \cite[\S 2--3]{Remondiere2026ND}.
\end{theorem}

\begin{remark}[The $2$-rank governs, $h_K$ does not]\label{rem:rk2-vs-hK}
Equation~\eqref{eq:univ-count} has the striking feature that the
$2$-rank $r = \mathrm{rk}_2\,\Cl(K)$ governs the orbit count, while the
class number $h_K = \lvert\Cl(K)\rvert$ does not appear explicitly. Two
anchors with the same $h_K$ but different $r$ produce different orbit
counts at any given weight; two anchors with different $h_K$ but
identical $(r, e)$ produce identical counts. The genus group
$\Cl(K)/\Cl(K)^2 \cong (\Z/2\Z)^r$ controls the \emph{cardinality} of the
$\Q$-rational tower; the exponent $\exp(\Cl(K))$ controls the
\emph{weight range} on which the tower is non-empty.
\end{remark}

\begin{example}[Verification across $r \in \{0,1,2,3,4\}$]\label{ex:hK-cyclic}
The empirical rule \eqref{eq:univ-count} has been verified by direct
\texttt{PARI/GP} computation at the following representative anchors:
\begin{itemize}
\item $h_K = 1$ (e.g.\ $D = -7$): $\Cl(K) = [1]$, $r = 0$, $e = 1$.
Predicted $1$ $\Q$-rat orbit at every odd $w \ge 3$. Confirmed.
\item $h_K = 2$ (e.g.\ $D = -15, -91$): $\Cl(K) = [2]$, $r = 1$,
$e = 2$. Predicted $2$ orbits at every odd $w \ge 3$. Confirmed
($664/664$ split primes).
\item $h_K = 4$ cyclic (e.g.\ $D = -56, -68$): $\Cl(K) \cong \Z/4\Z$,
$r = 1$, $e = 4$. Predicted $0$ orbits at $w = 3$ (since $4 \nmid 2$);
$2$ orbits at $w = 5, 9, 13, \ldots$ Confirmed at $w \in \{5, 9\}$.
\item $h_K = 4$ cyclic (e.g.\ $D = -39$): $\Cl(K) \cong \Z/4\Z$ with
$r = 1$, $e = 4$ (same structure as $D = -56, -68$).
\item $h_K = 4$ Klein (e.g.\ $D = -84, -120, -132$): $\Cl(K) \cong
(\Z/2\Z)^2$, $r = 2$, $e = 2$. Predicted $4$ orbits at every odd
$w \ge 3$. Confirmed at $D = -132$ ($96/96$).
\item $h_K = 8$ with $\Cl(K) \cong (\Z/2\Z)^3$ (e.g.\ $D = -420$):
$r = 3$, $e = 2$. Predicted $8$ orbits at every odd $w \ge 3$.
\item $h_K = 16$ with $\Cl(K) \cong (\Z/2\Z)^4$ (e.g.\ $D = -5460$):
$r = 4$, $e = 2$. Predicted $16$ orbits at every odd $w \ge 3$, in
agreement with a Faltings-style genus-rank sweep of fundamental
discriminants $|D| \le 20{,}000$ which finds $21$ anchors with $r = 4$
and zero with $r \ge 5$.
\end{itemize}
\end{example}

\begin{remark}[Status: TIER\_PROVED]\label{rem:thm-D-status}
Theorem~\ref{thm:univ-count} is \textit{TIER\_PROVED} at all $r \ge 1$
via the four-step chain established in the companion note
\cite[OP-THM-D-RIGOROUS]{Remondiere2026qD} :
(i) Hecke--Shimura--Ribet existence and parameterisation of
$\Q$-rational CM newforms; (ii) Sch\"utt 2008
\cite{Schutt2008} criterion for $\Q$-rationality requiring $\psi^{w-1}$
to factor through a genus character; (iii) explicit bijection
$\Q$-rat Hecke characters / $(\psi \leftrightarrow \bar\psi) \leftrightarrow
\widehat{\Cl(K)/\Cl(K)^2}$ of cardinality $2^r$; (iv) the converse
$e \nmid (w-1) \Rightarrow$ no $\Q$-rational Hecke character of
infinity type $(w-1, 0)$ exists, hence zero orbits.
Empirical support : $930/930$ split-prime Newton--Dickson identities
plus $171$ EXACT cross-Galois right-edge $\Q$-rationals at
$r \in \{0, 1, 2\}$, all PARI-verified at $50$--$80$-digit precision.
\end{remark}

\subsection{The $D = -91$ Hecke automorphism}\label{ss:open-D91}

\begin{conjecture}\label{conj:D91}
The structural identity
$\qstr{3}{11}(-91) = \qstr{7}{11}(-91) = 518/533$ of
Proposition~\ref{prop:D91-id} is induced by an explicit Hecke automorphism
of $S_w^{\mathrm{new}}(\Gamma_0(91), \chi_{-91})$ (for $w \in \{3, 7\}$)
arising from the splitting $\disc(K) = -91 = -7\cdot 13$ and the
isomorphism of genus fields
$K(\sqrt{-7}) \cong K(\sqrt{-13})$, which exchanges the weight-$3$ and
weight-$7$ towers via the Newton--Dickson iteration
\[
   a_p(f_7) = a_p(f_3)^3 - 3\,p\,a_p(f_3) \qquad (\text{Dickson}, D_3),
\]
under the involution $a_p(f_3)\mapsto \chi_{-13}(p)\,a_p(f_3)$
followed by the Dickson iteration to weight $11$.
\end{conjecture}

\subsection{Beilinson regulator class interpretation}\label{ss:open-Beilinson}

Each individual right-edge value $\alpha(\fl{w}{i}, w-1)$ is, after Beilinson
\cite{Beilinson1985} and Scholl \cite{Scholl1990}, the value of the
Beilinson regulator on an explicit class
\[
   \xi(\fl{w}{i}, w-1) \in \HM\bigl(\Spec K, M(\fl{w}{i})(w-1)\bigr)_\Q
\]
in motivic cohomology. The cross-Galois ratio $\qstr{w}{w'}(D)$ is the
corresponding ratio in the $\Q^{\times}$-quotient of the regulator image.
A fully motivic formulation, in the spirit of the Bloch--Kato Tamagawa
number conjecture for CM modular forms \cite{Castella2025}, would interpret
the empirical lattice of Theorem~\ref{thm:main-empirical} as a structure on
the cohomology group
\[
   \bigoplus_{(w, w')} \HM\bigl(\Spec K, M(\fl{w'}{1})(w'-1)\otimes
   M(\fl{w}{2})(w-1)\bigr)_\Q\Big/\Q^{\times}.
\]
We do not pursue this here.

\subsection{Cross-validation across class number}\label{ss:open-cross}

The Newton--Dickson chain
$a_p(f_9) = a_p(f_5)^2 - 2 p^4$ at weight pair $(5, 9)$ has been verified
to $332/332$ split primes at $h_K = 1$
\cite{Remondiere2026ND}; the same identity at $h_K = 2$ admits a per-orbit
signed form by Theorem~\ref{thm:newton-dickson}. The cross-orbit verification
at $h_K = 2$ for the wider Dickson family
$a_p(f_w) = D_{w-1}(X, p)$ with $X^2 = a_p(f_3) + 2p$
\cite[OP-DICKSON-UNIFIED]{Remondiere2026ND} has been verified at $120/120$
split primes at weight $w = 11$. A natural extension is to record the
$120 + 332 = 452$ uniform per-orbit signed identifications as a
\emph{single} arithmetic family across all odd weights $w \le 15$, with
the per-orbit twist $\chi_{d_i}^{(w'-w)/2}$ governing the sign at every
step. This is conjecturally a complete description of the Hecke action on
the $\Q$-rational CM tower at $h_K = 2$.

%===========================================================================
\section*{Acknowledgements}
%===========================================================================

The author thanks the LMFDB collaboration \cite{LMFDB} for the modular
forms database used throughout (forms of label
\texttt{N.w.a.\textbf{x}} for $\absD \in \{15, 20, 24, 35, 40, 51, 52, 88,
91, 132\}$, weights $w \in \{3, 5, 7, 9, 11, 13\}$), and the developers of
\texttt{PARI/GP} for the high-precision $L$-value evaluation routines
(\texttt{lfun}, \texttt{lfunmf}, \texttt{bestappr}, \texttt{lindep})
indispensable to this work.

%===========================================================================
\begin{thebibliography}{99}
%===========================================================================

\bibitem{Beilinson1985}
A.~A.~Beilinson,
\textit{Higher regulators and values of L-functions},
J.~Soviet Math.\ \textbf{30} (1985), 2036--2070.

\bibitem{BrunaultChida2015}
F.~Brunault and M.~Chida,
\textit{Regulators for Rankin--Selberg products of modular forms},
Ann.\ Math.\ Qu\'e.\ \textbf{40} (2016), 221--249.
\texttt{arXiv:1503.04626} [math.NT].

\bibitem{Castella2025}
F.~Castella,
\textit{Tamagawa number conjecture for CM modular forms and Rankin--Selberg
convolutions},
to appear in Proc.\ London Math.\ Soc., 2024 (revised 2025).
\texttt{arXiv:2407.11891} [math.NT].

\bibitem{ChowlaSelberg1949}
S.~Chowla and A.~Selberg,
\textit{On Epstein's zeta-function (I)},
Proc.\ Nat.\ Acad.\ Sci.\ U.S.A.\ \textbf{35} (1949), 371--374.

\bibitem{Cox1989}
D.~A.~Cox,
\emph{Primes of the form $x^2 + n y^2$: Fermat, class field theory, and
complex multiplication},
Wiley, 1989.
(Theorem 6.1 of Chapter 6 on genus characters and the bijection
$\Cl(K)/\Cl(K)^2\leftrightarrow\widehat{G_K^{\mathrm{gen}}}$ used here.)

\bibitem{Damerell1971a}
R.~M.~Damerell,
\textit{L-functions of elliptic curves with complex multiplication, I},
Acta Arith.\ \textbf{17} (1971), 287--301.

\bibitem{Gauss1801}
C.~F.~Gauss,
\emph{Disquisitiones Arithmeticae},
Leipzig, 1801.
(English transl.\ A.~A.~Clarke, Yale Univ.\ Press, 1965;
revised by W.~C.~Waterhouse, Springer, 1986.
Genus theory and prime-discriminant factorisation: \S\S 225--237.)

\bibitem{Gross1980}
B.~H.~Gross,
\emph{Arithmetic on elliptic curves with complex multiplication},
Lecture Notes in Math.\ \textbf{776}, Springer, 1980.

\bibitem{Hecke1936}
E.~Hecke,
\textit{\"Uber die Bestimmung Dirichletscher Reihen durch ihre
Funktionalgleichung},
Math.\ Ann.\ \textbf{112} (1936), 664--699.

\bibitem{KingsLoefflerZerbes2020}
G.~Kings, D.~Loeffler, and S.~L.~Zerbes,
\textit{Rankin--Eisenstein classes for modular forms},
Amer.\ J.\ Math.\ \textbf{142} (2020), 79--138.
\texttt{arXiv:1501.03289} [math.NT].

\bibitem{KingsSprang2025}
G.~Kings and J.~Sprang,
\textit{Eisenstein--Kronecker classes, integrality of critical values of
Hecke $L$-functions and $p$-adic interpolation},
Ann.\ Math.\ \textbf{202} (2025) (accepted; preprint 2019).
\texttt{arXiv:1912.03657} [math.NT].

\bibitem{KingsSprang2025survey}
G.~Kings and J.~Sprang,
\textit{Algebraicity of critical Hecke $L$-values},
to appear in Proc.\ Internat.\ Congress of Basic Science (ICBS) 2025, 2025.
\texttt{arXiv:2511.05198} [math.NT].

\bibitem{Kufner2024}
H.-U.~Kufner,
\textit{Deligne's conjecture on the critical values of Hecke $L$-functions},
preprint, 2024.
\texttt{arXiv:2406.06148} [math.NT].

\bibitem{Lidl1993}
R.~Lidl, G.~L.~Mullen, and G.~Turnwald,
\emph{Dickson polynomials},
Pitman Monographs and Surveys in Pure and Applied Math.\ \textbf{65},
Longman Scientific \& Technical, 1993.

\bibitem{LMFDB}
The LMFDB Collaboration,
\emph{The L-functions and Modular Forms Database},
\url{https://www.lmfdb.org}, accessed 2026.

\bibitem{Remondiere2026qD}
K.~Remondi\`ere,
\textit{A closed-form ratio identity for canonical pairs of weight-three
and weight-five CM newforms at class number two: Chowla--Selberg periods and
Beilinson--Damerell L-values},
companion note, 2026.

\bibitem{Remondiere2026ND}
K.~Remondi\`ere,
\textit{A Newton--Dickson recursion for Hecke eigenvalues of CM newforms
over imaginary quadratic fields},
companion note, 2026.

\bibitem{Scholl1990}
A.~J.~Scholl,
\textit{Motives for modular forms},
Invent.\ Math.\ \textbf{100} (1990), 419--430.

\bibitem{Shimura1971}
G.~Shimura,
\emph{Introduction to the Arithmetic Theory of Automorphic Functions},
Princeton Univ.\ Press, 1971.

\bibitem{Shimura1976}
G.~Shimura,
\textit{The special values of the zeta functions associated with cusp forms},
Comm.\ Pure Appl.\ Math.\ \textbf{29} (1976), 783--804.

\end{thebibliography}

\end{document}
